\begin{table}[t]
\centering
\caption{Klasik yöntemlerin yayımlanan koddaki uygulaması}
\begin{tabular}{lllll}
\toprule
Yöntem & Amaçta h\_d & Başlangıç & Güncelleme & Durma \\
\midrule
CD & var & bütün fazlar 0 & eleman başına kapalı form, sürekli faz & sabit K tarama (K = 5) \\
AO & var & bütün fazlar 0 & CD ile aynı güncelleme & sabit, K = 1 \\
SPA & var & h\_d & elemanları sırayla hizalama & tek geçiş \\
PGD & var & bütün fazlar 0 & gradyan adımı 0.2 + birim modüle izdüşüm & sabit 50 yineleme \\
"WMMSE" (betik) & var & bütün fazlar 0 & eleman başına bir gradyan adımı (0.1) & 1 yineleme \\
PSO & var (alınan güç) & rastgele, tekdüze & 20 parçacık & sabit 30 yineleme \\
Etiket (Denk. 8) & yok & — & --∠(G\_n F\_n) & — \\
\bottomrule
\end{tabular}
\end{table}
% Kaynak: aisy_paket/orijinal/ris_CD.m; ris_AO/SPA/PGD/WMMSE/PSO.m satır satır aktarımı audit/scripts/step2_baselines.py
